\subsection{What the trained fields learned}
\label{sec:numerics}

Trajectory error measures one orbit. Because we saved the checkpoints, we
can also examine the learned field $\mathbf{f}_\theta$ itself as a dynamical
system, using the invariants of Section~\ref{sec:datasets} as ground truth.

\paragraph{Equilibria by Newton's method.} We solve
$\mathbf{f}_\theta(\mathbf{u}^\ast)=\mathbf{0}$ by Newton-Raphson,
$\mathbf{u}\leftarrow\mathbf{u}-J_\theta(\mathbf{u})^{-1}\mathbf{f}_\theta(\mathbf{u})$,
with the Jacobian from automatic differentiation in float64, starting at the
true equilibrium $(3,1.5)$. The iteration converges quadratically for every
working model. For the Tsit5 model the residual norms are
$2.1,\ 1.6\times10^{-1},\ 1.1\times10^{-2},\ 9.5\times10^{-6},\ 3.4\times10^{-12}$,
then exactly zero. The residual is roughly squared at every step, which is
the signature of quadratic convergence ($\|\mathbf{F}_{k+1}\|/\|\mathbf{F}_k\|^2$
stays between $0.04$ and $0.43$). The roots, however, are far from the truth
(Table~\ref{tab:equilibria}). Every RBF solver model places its equilibrium
near $(1.9,1.04)$, about $1.15$ from $(3,1.5)$, and it is an unstable focus
with $\mathrm{Re}\,\lambda\approx+1.2$, where the true system has
$\mathrm{Re}\,\lambda=0$. The imaginary parts are close to the true $2.12$.
The Chebyshev and Lagrange models have \emph{stable} foci instead. Newton
iteration on the Chebyshev model does not converge within 30 steps, and on
the Newton-basis model it diverges.

\begin{table*}[t]
\centering
\footnotesize
\renewcommand{\arraystretch}{1.08}
\caption{Numerical analysis of the trained fields. $\mathbf{u}^\ast$:
learned equilibrium (Newton-Raphson from $(3,1.5)$); $\lambda$: eigenvalues
of $J_\theta(\mathbf{u}^\ast)$; period error from zero-crossings of the
predicted prey signal over $[0,28]$; $H$ drift is
$\max_t|H(\hat{\mathbf{u}})-H_0|/H_0$ on $[0,28]$. True values:
$\mathbf{u}^\ast=(3,1.5)$, $\lambda=\pm2.121i$, $T=3.3175$, drift $0$.}
\label{tab:equilibria}
\begin{tabular}{@{}L{4.0cm} C{2.6cm} C{2.0cm} C{1.6cm} C{1.6cm} C{2.0cm} C{1.4cm}@{}}
\toprule
\hdrrow \thh{Model} & \thh{$\mathbf{u}^\ast$} & \thh{$\|\mathbf{u}^\ast-\mathbf{u}^\ast_{\text{true}}\|$} & \thh{$\mathrm{Re}\,\lambda$} & \thh{$|\mathrm{Im}\,\lambda|$} & \thh{Period error} & \thh{$H$ drift}\\
\midrule
Euler & $(1.80,\,1.02)$ & $1.30$ & $+0.51$ & $2.05$ & $-1.7\times10^{-3}$ & $4.2\%$\\
Heun & $(1.91,\,1.03)$ & $1.19$ & $+1.11$ & $2.03$ & $-4.7\times10^{-4}$ & $3.1\%$\\
Midpoint & $(1.93,\,1.04)$ & $1.16$ & $+1.17$ & $2.01$ & $\mathbf{-3.5\times10^{-6}}$ & $3.1\%$\\
RK4 & $(1.94,\,1.05)$ & $1.15$ & $+1.21$ & $1.99$ & $+7.8\times10^{-5}$ & $3.2\%$\\
DOPRI5 & $(1.94,\,1.04)$ & $1.16$ & $+1.20$ & $2.00$ & $+3.5\times10^{-5}$ & $3.2\%$\\
Tsit5 / RBF & $(1.94,\,1.04)$ & $1.15$ & $+1.20$ & $2.01$ & $+1.7\times10^{-5}$ & $3.1\%$\\
\midrule
B-spline & $(1.66,\,0.94)$ & $1.46$ & $+1.18$ & $2.99$ & $+1.3\times10^{-4}$ & $\mathbf{1.3\%}$\\
Lagrange & $(2.67,\,2.15)$ & $0.73$ & $-0.44$ & $5.21$ & $-9.9\times10^{-4}$ & $3.9\%$\\
Chebyshev & \multicolumn{4}{c}{Newton iteration does not converge} & $+6.0\times10^{-4}$ & $5.1\%$\\
IQF & $(2.05,\,1.19)$ & $1.00$ & $+0.26$ & $1.70$ & $-2.2\times10^{-3}$ & $76\%$\\
RSWAF & $(2.15,\,0.91)$ & $1.03$ & $+1.00$ & $1.69$ & $+6.3\times10^{-4}$ & $5.4\%$\\
Newton & \multicolumn{4}{c}{Newton iteration diverges} & $-6.4\times10^{-3}$ & $12\%$\\
\midrule
MLP-ODE (SiLU), $10^4$ & $(2.35,\,1.12)$ & $0.76$ & $+0.30$ & $2.75$ & $+8.0\times10^{-4}$ & $2.8\%$\\
KAN-ODE, $5\times10^4$ & $(1.51,\,0.80)$ & $1.65$ & $+0.67$ & $2.85$ & $-1.6\times10^{-4}$ & $0.9\%$\\
MLP-ODE (SiLU), $5\times10^4$ & $(3.65,\,1.30)$ & $\mathbf{0.68}$ & $+0.95$ & $1.83$ & $+1.0\times10^{-4}$ & $\mathbf{0.4\%}$\\
\bottomrule
\end{tabular}
\end{table*}

\paragraph{The learned system is a limit-cycle oscillator.} An unstable
focus surrounded by an orbit that trajectories stay on is the signature of an
attracting limit cycle, not of a conservative centre. We confirm this by
integrating each learned field from four initial conditions whose true orbits
have $H\in\{0.85,\,0.86,\,2.0,\,3.5\}$ (Figure~\ref{fig:limit-cycle}). In the
true system the four orbits are distinct nested curves. Under every learned
field (KAN and MLP, at $10^4$ and at $5\times10^4$ epochs) all four
trajectories converge onto the single training orbit, with late-time $H$
between $2.000$ and $2.008$ for every starting point. Both architectures
have replaced Lotka-Volterra's one-parameter family of neutral orbits with
one isolated, attracting cycle through the training data. This explains two
things at once. First, the flat extrapolation error of the KAN-ODE
(Table~\ref{tab:extrap-horizon}) is partly a property of the attractor:
amplitude errors are damped back onto the cycle, so only phase error can
accumulate. That is why the period error in Table~\ref{tab:equilibria}
predicts the extrapolation ranking better than any field metric. Second, the
models say nothing reliable about any other initial condition. They have
learned an orbit, not the vector field.

\begin{figure*}[t]
\centering
\includegraphics[width=\textwidth]{figures/analysis/limit_cycle_test.pdf}
\caption{Trajectories of each learned field from four initial conditions
(coloured dots) over $t\in[0,40]$, against the four corresponding true orbits
(grey). Red crosses mark the learned equilibria, black plus signs the true
one. Every learned field attracts all four onto the training orbit.}
\label{fig:limit-cycle}
\end{figure*}

\paragraph{Where the field is right.} Figure~\ref{fig:field-error} maps the
field error $\|\mathbf{f}_\theta-\mathbf{f}\|$ over the phase plane. It is
small only in a thin tube around the training orbit. On the orbit the error
is $2.7$--$2.8\%$ of the mean field magnitude for both 10k models and
$0.6$--$1.1\%$ at 50k. Across the plotted box the median error is
$20$--$34\%$, and for every model less than $3\%$ of the box is below $1\%$
error. Longer training sharpens the tube without widening it: the 50k KAN's
median off-orbit error ($34\%$) is higher than the 10k model's ($24\%$).

\begin{figure*}[t]
\centering
\includegraphics[width=0.9\textwidth]{figures/analysis/field_error_maps.pdf}
\caption{Field error $\|\mathbf{f}_\theta(\mathbf{u})-\mathbf{f}(\mathbf{u})\|$,
normalised by the mean true field magnitude on the orbit ($4.82$), clipped to
$[10^{-3},1]$. Thick black: training segment; blue: rest of the orbit;
plus sign: true equilibrium.}
\label{fig:field-error}
\end{figure*}

\paragraph{Invariant and error growth.} Figure~\ref{fig:first-integral}b
tracks the first integral along each solver model's own prediction. All of
them hold $H$ to within $3$--$4\%$, and the drift is oscillatory rather than
secular, as it must be on an attracting cycle. The learned models' invariant
error is dominated by model error, orders of magnitude above the solver
drift in panel (a), except for Euler. Figure~\ref{fig:error-vs-time} shows the
pointwise error out to $t=28$. The error envelopes of the Midpoint and Tsit5
models grow by $5.5$--$5.7\times10^{-4}$ per unit time, those of Heun and
Chebyshev by $6$--$8\times10^{-3}$, and Euler's by $2.1\times10^{-2}$.
Panel (b) previews Section~\ref{sec:phase4}. At $5\times10^4$ epochs the MLP
error is lower than the KAN's throughout, while the 50k KAN has a
\emph{larger} far-horizon error than the 10k KAN (maximum $6.5\times10^{-2}$
against $3.4\times10^{-2}$ on $(14,28]$), despite its better scored-window
MSE.

\begin{figure*}[t]
\centering
\includegraphics[width=\textwidth]{figures/analysis/first_integral_drift.pdf}
\caption{Relative drift of the Lotka-Volterra first integral $H$
(Equation~\ref{eq:lv-invariant}). (a) The exact field integrated by each
solver at $h=0.05$: pure truncation drift (Euler's trajectory leaves the
domain of $H$ at $t=19.4$). (b) Each trained model integrated by its
training solver. Shaded: training window; dotted: end of the scored horizon.}
\label{fig:first-integral}
\end{figure*}

\begin{figure*}[t]
\centering
\includegraphics[width=0.9\textwidth]{figures/analysis/error_vs_time.pdf}
\caption{Pointwise trajectory error on $[0,28]$. (a) Solver and basis
choice at $10^4$ epochs (Midpoint and Tsit5 overlap). (b) KAN-ODE and
MLP-ODE at $10^4$ and $5\times10^4$ epochs.}
\label{fig:error-vs-time}
\end{figure*}

\paragraph{What the edges look like.} Figure~\ref{fig:edges} plots the
twenty first-layer edge functions $\phi_{i,j}$ of the RBF and B-spline
models over the range each input covers on the orbit. Over that range most
edges are close to affine, with mild curvature concentrated at small
populations, where $\tanh$ is least saturated. None of them resembles a term
of Equation~\ref{eq:lv}. The bilinear products $u_1u_2$ are assembled by
composition through the second layer rather than carried by any single
edge, which is consistent with the pruning result of
Section~\ref{sec:track-e}.

\begin{figure*}[t]
\centering
\includegraphics[width=\textwidth]{figures/analysis/edge_functions.pdf}
\caption{First-layer edge functions (spline plus residual branch) of the
trained RBF and B-spline models, each input varied over its range on the
orbit with the other held at its mean. The edge with the largest range is
highlighted.}
\label{fig:edges}
\end{figure*}

\subsubsection{Training dynamics across the sweeps}

Figures~\ref{fig:ablation-curves} and~\ref{fig:basis-curves} show the test
MSE of every configuration during training, and Table~\ref{tab:milestones}
records when each one's training MSE crosses fixed thresholds. Five of the
six solver curves lie on top of each other throughout. Euler separates from
about epoch 2000 and stays above the others until the final approach, so
solver choice changes where training ends rather than how it starts. From
about epoch 6000 every solver shows periodic loss spikes of one to two
orders of magnitude that recover within a few hundred epochs. The basis
curves separate much earlier. B-spline is lowest almost throughout and Newton
highest, and B-spline, Chebyshev and RSWAF reach $10^{-2}$ training MSE
first ($1900$--$2050$ epochs, against RBF's $3112$). Only B-spline turns
that head start into better extrapolation.

\begin{table*}[t]
\centering
\footnotesize
\renewcommand{\arraystretch}{1.08}
\caption{First epoch at which the training MSE falls below each threshold
(``--'': not reached within $10^4$ epochs).}
\label{tab:milestones}
\begin{tabular}{@{}L{2.8cm} C{1.4cm} C{1.4cm} C{1.4cm} @{\hspace{1.2cm}} L{2.8cm} C{1.4cm} C{1.4cm} C{1.4cm}@{}}
\toprule
\hdrrow \thh{Solver} & \thh{$10^{-2}$} & \thh{$10^{-3}$} & \thh{$10^{-4}$} & \thh{Basis} & \thh{$10^{-2}$} & \thh{$10^{-3}$} & \thh{$10^{-4}$}\\
\midrule
Euler & 3533 & 6729 & -- & B-spline & \textbf{1894} & 4826 & \textbf{8458}\\
Heun & 3144 & 5564 & 9512 & Gaussian RBF & 3112 & 5282 & 9617\\
Midpoint & \textbf{3102} & 5285 & 9599 & Chebyshev & 1931 & \textbf{3892} & 9828\\
RK4 & 3153 & 5318 & 9798 & Lagrange & 2924 & 5407 & --\\
DOPRI5 & 3132 & 5305 & 9796 & IQF & 2561 & 5726 & --\\
Tsit5 & 3112 & \textbf{5282} & 9617 & RSWAF & 2047 & 4177 & --\\
MLP-ODE (SiLU) & 741 & 1297 & 9145 & Newton & 4577 & -- & --\\
\bottomrule
\end{tabular}
\end{table*}

\begin{figure}[t]
\centering
\includegraphics[width=\columnwidth]{figures/ablation/05_solver_convergence.png}
\caption{Test MSE during training for the six solvers (RBF basis),
logged every ten epochs; the dotted line is $10^{-4}$.}
\label{fig:ablation-curves}
\end{figure}

\begin{figure}[t]
\centering
\includegraphics[width=\columnwidth]{figures/ablation/06_basis_convergence.png}
\caption{Test MSE during training for the seven bases (Tsit5 solver),
logged every ten epochs; the dotted line is $10^{-4}$.}
\label{fig:basis-curves}
\end{figure}

\begin{keypoint}[What the ablation and the numerical analysis establish]
\begin{itemize}
\item All six integrators satisfy their order conditions to round-off and
  show their design order; Tsit5's leading error constant is $2.9\times$
  smaller than DOPRI5's.
\item Above order two, solver choice does not change accuracy at $h=0.05$;
  Midpoint is the cost optimum (Tsit5's accuracy at $29\%$ of the time).
\item The network learns the field that its training solver needs, not the
  physical field: the Euler model is an accurate model of the inverse
  modified equation, and models do not transfer across solver families.
\item Locality, not conditioning, separates the bases; on this system a third
  of the first layer's grid is idle because states are positive.
\item Every trained model, KAN or MLP, is an attracting limit-cycle
  oscillator around a spurious unstable equilibrium. It extrapolates the
  training orbit well and says nothing reliable about any other.
\end{itemize}
\end{keypoint}
